% R18_Reclassification_Audit.tex
% monogate Cost Theory — R18: Reclassification Audit
% Verdicts on T35, T36, T37, T40 (currently labelled PROPOSITION)
% Date: 2026-04-20
% Author: Arturo R. Almaguer

\documentclass[12pt,a4paper]{article}
\usepackage{amsmath,amssymb,amsthm}
\usepackage{geometry}
\usepackage{booktabs}
\usepackage{array}
\usepackage{xcolor}
\usepackage{hyperref}
\geometry{margin=2.5cm}

% ── Theorem environments ─────────────────────────────────────────────────────
\newtheorem{theorem}{Theorem}[section]
\newtheorem{lemma}[theorem]{Lemma}
\newtheorem{corollary}[theorem]{Corollary}
\newtheorem{proposition}[theorem]{Proposition}
\theoremstyle{definition}
\newtheorem{definition}[theorem]{Definition}
\newtheorem{remark}[theorem]{Remark}

% ── Macros ───────────────────────────────────────────────────────────────────
\newcommand{\Cost}{\operatorname{Cost}}
\newcommand{\NC}{\operatorname{NaiveCost}}
\newcommand{\SB}{\operatorname{SB}}
\newcommand{\ceil}[1]{\left\lceil #1 \right\rceil}
\newcommand{\Fam}{\mathcal{F}}

% ── Verdict colour boxes ──────────────────────────────────────────────────────
\newcommand{\verdictTHEOREM}{\textbf{\textcolor{blue}{VERDICT: THEOREM}}}
\newcommand{\verdictPROP}{\textbf{\textcolor{red}{VERDICT: PROPOSITION}}}

\title{Reclassification Audit (R18)\\[6pt]
       \large T35, T36, T37 (Lower Bounds) and T40 (Linear Cost Law)\\
       Should They Be THEOREM or PROPOSITION?}
\author{Arturo R.\ Almaguer\\
\small monogate research \quad \texttt{almaguer1986@gmail.com}}
\date{April 2026}

% ═════════════════════════════════════════════════════════════════════════════
\begin{document}
\maketitle

% ── Preamble ─────────────────────────────────────────────────────────────────
\begin{abstract}
This audit reviews four results currently labelled \textsc{Proposition}
in the monogate cost theory series and determines whether each qualifies
for reclassification as a \textsc{Theorem}.  The guiding criterion is:
\emph{``being minor does not make something a Proposition --- being
straightforward does.''}  A result with a complete, gap-free proof,
however brief, is a Theorem if it establishes a non-trivial fact via
structured argument rather than unfolding a definition verbatim.

The four results examined are:
\begin{itemize}
  \item \textbf{T35} --- Cost$(E) \geq 0$ (non-negativity).
  \item \textbf{T36} --- Cost$(E) \geq \lceil\mathrm{depth}(E)/2\rceil$
        (depth lower bound).
  \item \textbf{T37} --- Cost$(E) \geq n_{\exp} + n_{\ln}$
        (operation-type lower bound).
  \item \textbf{T40} --- For $N$-term positive-domain summations:
        $\Cost = (\alpha_0 + 3) \cdot N - 3$ exactly
        (the Summation Scaling Formula; also identified as T40b in R5
        and as T42 in the complete theory index).
\end{itemize}

\noindent\textbf{Summary of verdicts:} T35 THEOREM, T36 THEOREM, T37 THEOREM,
T40 (Summation Scaling Formula) THEOREM.
\end{abstract}

\tableofcontents
\bigskip

% ═════════════════════════════════════════════════════════════════════════════
\section{Audit Criteria}
% ═════════════════════════════════════════════════════════════════════════════

\begin{definition}[Theorem vs.\ Proposition standard]
A result is classified as a \textbf{Theorem} when its proof:
\begin{enumerate}
  \item is complete and gap-free;
  \item involves a structural argument beyond verbatim definition-unfolding;
  \item would require non-trivial effort to reconstruct without the proof.
\end{enumerate}
A result is classified as a \textbf{Proposition} when its proof reduces
entirely to substituting a definition, or when a critical step is missing,
assumed without argument, or circular.  Brevity alone is not grounds for
demotion: a two-line proof that closes a non-trivial inference is a Theorem.
\end{definition}

% ═════════════════════════════════════════════════════════════════════════════
\section{T35: Non-negativity of Cost}
% ═════════════════════════════════════════════════════════════════════════════

\subsection{Statement}

The statement appearing in the literature has two distinct formulations:

\begin{enumerate}
  \item \textbf{R1 version (P1):}  For every expression tree
        $E \in \mathcal{E}(\mathcal{T},\Fam_{16})$,
        \[
          \Cost(E) \;\geq\; 0.
        \]
  \item \textbf{COST-3 version (T35, Trivial Lower Bound):}
        For any EML-family expression $E$ computing a target function $f$,
        \[
          \Cost(E) \;\geq\; |\mathcal{O}(E)|,
        \]
        where $\mathcal{O}(E)$ is the set of distinct non-mergeable operator
        instances in $E$.  The Cost\_Theory\_Complete.tex index summarises this
        (conservatively) as $\Cost(E) \geq 1$ for any non-terminal $E$.
\end{enumerate}

The P1/non-negativity variant is purely definitional.  The T35/trivial-lower-bound
variant is structural.  We audit both.

\subsection{Proof Analysis}

\subsubsection{P1 (Non-negativity, $\Cost(E) \geq 0$)}

\textbf{Proof in R1 (verbatim key steps):}
By Definition of Cost, $\Cost(E)$ is the minimum of the set
$S = \{|V_{\mathrm{int}}(D)|\}$ over all valid DAGs $D$.  Every element of
$S$ is a cardinality, hence $\geq 0$.  Therefore $\min S \geq 0$.

\textbf{Assessment:}  This is a one-line consequence of the fact that
$\Cost$ is defined as a minimum of a set of natural numbers.  In Lean 4 the
proof reduces to \texttt{Nat.zero\_le \_}: the type system makes it
vacuous.  The proof does not establish any non-trivial structural fact;
it merely confirms that the definition lands in $\mathbb{N}_{\geq 0}$.
This is the paradigmatic case of a definitional consequence.

\textbf{Gap analysis:}  No gaps.  The argument is complete.  But it is
also entirely definitional: substitute the definition, observe the range
is $\mathbb{N}_0$, conclude.

\textbf{Verdict for P1 alone:}
Were P1 the only content, it would remain a \textbf{Proposition} (purely
definitional).  However, T35 in COST-3 subsumes P1 and establishes a
strictly stronger statement.

\subsubsection{T35 (Trivial Lower Bound, $\Cost(E) \geq |\mathcal{O}(E)|$)}

\textbf{Proof in COST-3 (key steps):}
\begin{enumerate}
  \item Each element of $\mathcal{O}(E)$ requires at least one distinct
        operator node in any valid $\Fam_{16}$ tree.
  \item Elements of $\mathcal{O}(E)$ are by definition non-mergeable, so no
        two elements share a single node.
  \item Therefore the total node count $\geq |\mathcal{O}(E)|$.
\end{enumerate}

\textbf{Assessment:}  The critical step (2) is a structural argument: it
appeals to the definition of non-mergeability in $\Fam_{16}$ to establish
that distinct operations require distinct nodes.  This is not merely
unwrapping the definition of $\Cost$; it requires knowing that the 16-operator
family does \emph{not} allow arbitrary multi-operation fusion.  The proof is
2--3 lines and complete.

The remark on merging (COST-3, Remark on $\Fam_{16}$ operators) correctly
notes that failing to account for merging gives an over-estimate, so the
subtlety of what counts as $|\mathcal{O}|$ is genuine.

\textbf{Gap analysis:}  None.  The proof is complete and the structural step
is sound.  The worked examples (pH, $\Delta G$) confirm the bound is
non-trivially tight or non-vacuous.

\textbf{Overall verdict for T35:}  \verdictTHEOREM.
The non-mergeability argument is structural, not definitional.  The bound
$\Cost(E) \geq |\mathcal{O}(E)|$ is a genuine lower bound theorem.

% ═════════════════════════════════════════════════════════════════════════════
\section{T36: Depth Lower Bound}
% ═════════════════════════════════════════════════════════════════════════════

\subsection{Statement (COST-3)}

Let $f$ be a target function and suppose any correct EML-family computation of
$f$ must produce at least $d$ distinct intermediate values.  Then:
\[
  \SB(f) \;\geq\; \left\lceil \frac{d}{2} \right\rceil.
\]
The proof also establishes the sharper bound $\SB(f) \geq d$ (one intermediate
per node) and notes the $\lceil d/2 \rceil$ form is the depth-layer variant.

\subsection{Proof Analysis}

\textbf{Key steps in COST-3:}
\begin{enumerate}
  \item An operator node $\op(E_1, E_2)$ produces exactly one new intermediate
        value (its own output).
  \item A tree of $k$ nodes contains exactly $k$ internal nodes and $k+1$
        leaves (the standard binary-tree identity).  Each internal node
        produces one intermediate.  Hence $k$ nodes produce $\leq k$
        intermediates, so $k \geq d$: thus $\SB(f) \geq d$.
  \item The stated $\lceil d/2 \rceil$ is the weaker depth-layer bound: at each
        depth layer, at most 2 intermediates are produced per node (two children
        can both be intermediate).  Taking the ceiling over depth layers gives
        $\lceil d/2 \rceil$.
\end{enumerate}

\textbf{Assessment:}  The sharper bound $\SB(f) \geq d$ is a clean structural
theorem using the binary-tree internal-node count identity, which itself
requires a (standard, short) inductive proof.  The key fact ``each internal
node produces exactly one intermediate'' is not a definition; it follows from
the semantics of operator application.

The $\lceil d/2 \rceil$ formulation is weaker than $d$ and relies on a
depth-layer argument.  The proof correctly acknowledges this: ``the sharper
bound is in fact $\SB(f) \geq d$''.  The depth-layer version is included for
comparison with depth-based arguments elsewhere.

\textbf{Gap analysis:}  There is a mild presentational gap: the proof
conflates the sharper bound ($\geq d$) with the stated bound ($\geq \lceil d/2
\rceil$), offering the latter as the headline statement even though the
argument immediately yields the former.  This is not a logical gap but a
notational choice that could mislead.  Importantly, the weaker $\lceil d/2
\rceil$ bound is provably correct and follows as a corollary of $\geq d$.
There is no error; the statement is conservative.

\textbf{Verdict for T36:}  \verdictTHEOREM.
The proof establishes a structural lower bound via the binary-tree node-count
identity.  The argument is complete and involves a genuine structural
observation (one intermediate per node, binary tree identity).  The Corollary
explicitly states the tighter $\SB(f) \geq d$ form.

% ═════════════════════════════════════════════════════════════════════════════
\section{T37: Operation-Type Lower Bound}
% ═════════════════════════════════════════════════════════════════════════════

\subsection{Statement (COST-3)}

Let $E$ be an expression.  Denote by $n_{\exp}(E)$ the number of
non-constant exponential occurrences and by $n_{\ln}(E)$ the number of
non-constant logarithm occurrences.  Then:
\[
  \Cost(E) \;\geq\; n_{\exp}(E) \cdot 1 \;+\; n_{\ln}(E) \cdot 1.
\]

\subsection{Proof Analysis}

\textbf{Key steps in COST-3:}
\begin{enumerate}
  \item \emph{Exp lower bound.}  Each occurrence of $e^{g(x)}$ where $g(x)$
        is not constant must appear as the output of some node in the
        $\Fam_{16}$ tree.  A single $\Fam_{16}$ node computes a fixed binary
        function; distinct exponential occurrences (with distinct non-constant
        arguments) cannot be collapsed into one node.  Each contributes $\geq 1$
        to $\Cost(E)$.
  \item \emph{Ln lower bound.}  Identical argument for $\ln$.
  \item \emph{Additivity.}  Exp-nodes and ln-nodes are distinct in $\Fam_{16}$
        (each operator has a fixed functional form); they cannot share a node.
        Therefore the two contributions are additive.
\end{enumerate}

\textbf{Assessment:}  This is a structural counting argument.  The critical
step is that distinct non-constant exponential (or logarithmic) occurrences
cannot be collapsed: this uses the structure of $\Fam_{16}$, specifically that
no single binary operator simultaneously computes two independent transcendental
evaluations.  This is not merely a definition-unwrap.

The additivity step (exp-nodes and ln-nodes are disjoint) follows from
inspecting the $\Fam_{16}$ operator table: every operator has one functional
role, so an exp-node cannot simultaneously serve as an ln-node.  This is
structural rather than definitional.

\textbf{Gap analysis:}  The proof is complete.  Remark~(T37 tightness) correctly
notes the bound is tight for pure transcendental expressions and identifies
where it is not tight (polynomial/rational auxiliary operations).  The gap
analysis is honest and does not attempt to overstate the bound.

One minor imprecision: the proof says ``every operator node that produces
$e^{(\cdot)}$ in its output requires at least 1 dedicated node.''  This is
true but could be stated more precisely: in $\Fam_{16}$, every operator that
outputs a term involving $e^{g(x)}$ for non-constant $g$ is distinct from
operators outputting $e^{h(x)}$ for $h \neq g$, because the operator's binary
input structure fixes its output.  This is standard and sound.

\textbf{Verdict for T37:}  \verdictTHEOREM.
The argument is a counting argument over the structure of $\Fam_{16}$, not a
definition unwrap.  The proof is 3--4 sentences, complete, and correct.

% ═════════════════════════════════════════════════════════════════════════════
\section{T40: Summation Scaling Formula}
% ═════════════════════════════════════════════════════════════════════════════

\subsection{Identification and Statement}

In R5\_Linear\_Cost\_Law.tex the statement appears as Theorem T40b
(Summation Scaling Formula).  In the Cost\_Theory\_Complete.tex index it
appears as T42 (O(N) Scaling Law).  The task prompt refers to this as T40.
All three labels point to the same result:

\medskip
\noindent\textbf{Statement.}  Let $E_N = \sum_{i=1}^{N} f(i)$ where each
term $f(i)$ has fixed na\"{i}ve cost $\alpha_0$ per term, there are no shared
subexpressions across distinct terms, and no compound patterns apply.
For positive-domain summation ($c_{\mathrm{add}} = 3$):
\[
  \Cost(E_N) \;=\; (\alpha_0 + 3)\,N \;-\; 3 \quad \text{exactly}.
\]

\subsection{Base Case: $N = 1$}

For $N = 1$, the formula gives:
\[
  \Cost(E_1) = (\alpha_0 + 3) \cdot 1 - 3 = \alpha_0.
\]
A single-term expression has no \texttt{add} nodes: the cost is simply the
per-term cost $\alpha_0 = \NC(f(1))$.  Since there is no sharing and no
compound pattern by hypothesis, T38 gives $\Cost(E_1) = \NC(E_1) = \alpha_0$.
\textbf{Base case established.}

\subsection{Inductive Step}

\textbf{Claim:}  Adding one term from $E_N$ to $E_{N+1}$ increases cost
by exactly $\alpha_0 + 3$.

\textbf{Argument:}
\begin{enumerate}
  \item The new term $f(N+1)$ is an independent sub-expression with per-term
        cost $\alpha_0$ (no sharing with prior terms by hypothesis; no compound
        patterns by hypothesis).  By T38: $\Cost(f(N+1)) = \alpha_0$.
  \item Combining $E_N$ with $f(N+1)$ requires one binary \texttt{add} node
        (positive-domain variant, cost 3).  By the Additive Cost Law (T40-ADD):
        \[
          \Cost(E_{N+1}) = \Cost(E_N) + \Cost(f(N+1)) + c_{\mathrm{add}}
                         = \Cost(E_N) + \alpha_0 + 3.
        \]
  \item By the inductive hypothesis $\Cost(E_N) = (\alpha_0+3)N - 3$, so:
        \[
          \Cost(E_{N+1}) = (\alpha_0+3)N - 3 + \alpha_0 + 3
                         = (\alpha_0+3)(N+1) - 3. \quad \checkmark
        \]
\end{enumerate}

\textbf{Closed-form derivation (direct):}
The $N$ per-term sub-trees each cost $\alpha_0$, contributing $N\alpha_0$.
The $(N-1)$ binary \texttt{add} nodes joining the partial sums each cost 3,
contributing $3(N-1)$.  Total:
\[
  \NC(E_N) = N\alpha_0 + 3(N-1) = (\alpha_0 + 3)N - 3.
\]
Under no-sharing, no-pattern hypothesis: $\Cost(E_N) = \NC(E_N)$ by T38.

\subsection{Role of the Positive-Domain Assumption}

The assumption that all terms have positive domain means $c_{\mathrm{add}} = 3$
rather than 11.  If this assumption were dropped, the formula becomes
$(\alpha_0 + 11)N - 11$.  The formula as stated is specific to
$c_{\mathrm{add}} = 3$ and this specialisation is a genuine hypothesis,
not vacuous.

The positive-domain property is preserved compositionally: $e^u > 0$ always,
and a sum of positive terms is positive, so the cost assignment is consistent
throughout the tree.

\subsection{Gap Analysis}

\textbf{Potential gaps examined:}

\begin{enumerate}
  \item \textbf{Independence of terms.}  The proof assumes terms $f(i)$ are
        independent (no cross-term sharing).  This is explicitly stated and
        is enforced by the hypothesis ``no shared subexpressions across terms.''
        The most natural case is when each $f(i)$ depends on a distinct index
        variable $x_i$; the proof notes this.  No gap here.

  \item \textbf{Inductive cost of the add node.}  The claim is that each
        \texttt{add} node costs exactly 3.  This follows from the SuperBEST v3
        routing table (positive-domain add = 3) and the domain-folding theorem
        T40-DF, which justifies using cost 3 when all operands are
        statically verifiable as positive.  The proof assumes this
        qualification holds for the entire sum tree.  No gap for the stated
        hypothesis.

  \item \textbf{The Additive Cost Law (T40-ADD) as a dependency.}  The
        inductive step invokes T40-ADD (independent branches sum exactly).
        T40-ADD is itself classified as a \textsc{Proposition} in the
        complete theory index, but it is separately proved and its proof
        is sound.  Its use here is valid.

  \item \textbf{Exact equality vs.\ inequality.}  The formula gives an
        \emph{exact} cost, not a bound.  This exactness requires T38
        (Cost Decomposition) to ensure $\SD = \PB = 0$ under the stated
        hypotheses.  The proof invokes T38 for this step.  This is correct:
        under no-sharing and no-pattern, both reductions vanish by T38.

  \item \textbf{The $-3$ constant.}  The formula has $-3$ (not $-c_{\mathrm{add}}$
        in general).  For positive-domain this is $-3$ since $c_{\mathrm{add}} = 3$.
        This is consistent with the closed-form derivation: $(N-1)$ add nodes
        contribute $3(N-1) = 3N - 3$, and $N\alpha_0 + 3N - 3 = (\alpha_0+3)N - 3$.
        Correct.
\end{enumerate}

No gaps found.  The proof is complete under the stated hypotheses.

\subsection{Verdict for T40}

\verdictTHEOREM.

The proof is by induction on $N$ (or equivalently by direct closed-form
counting), with the inductive step established in 2--3 lines via T40-ADD
and T38.  The ``positive-domain'' assumption is a genuine hypothesis that
limits the scope of the result to a meaningful special case (all four
concrete instances in the literature are of this type).  The proof is
complete, the base case is verified, the inductive step is sound, and the
exact formula is established.

% ═════════════════════════════════════════════════════════════════════════════
\section{Summary and Reclassification Table}
% ═════════════════════════════════════════════════════════════════════════════

\begin{table}[h]
\centering
\caption{Reclassification audit results for T35, T36, T37, T40.}
\label{tab:audit}
\smallskip
\begin{tabular}{@{}clp{5.5cm}cc@{}}
\toprule
\textbf{Label} & \textbf{Statement} & \textbf{Proof character}
  & \textbf{Gap?} & \textbf{Verdict} \\
\midrule
T35
  & $\Cost(E) \geq |\mathcal{O}(E)|$
  & Non-mergeability structural argument over $\Fam_{16}$
  & None
  & \textcolor{blue}{\textbf{THEOREM}} \\[6pt]
T36
  & $\SB(f) \geq \lceil d/2 \rceil$
    (and $\SB(f) \geq d$ sharper)
  & Binary-tree internal-node count identity; one-intermediate-per-node
  & None (conservative statement only)
  & \textcolor{blue}{\textbf{THEOREM}} \\[6pt]
T37
  & $\Cost(E) \geq n_{\exp} + n_{\ln}$
  & Counting argument: distinct non-constant transcendental occurrences require
    distinct nodes in $\Fam_{16}$; additivity from operator distinctness
  & None
  & \textcolor{blue}{\textbf{THEOREM}} \\[6pt]
T40
  & $\Cost(E_N) = (\alpha_0+3)N - 3$ (positive domain)
  & Induction on $N$; base case $N=1$ direct; step via T40-ADD + T38;
    closed-form derivation also given
  & None
  & \textcolor{blue}{\textbf{THEOREM}} \\
\bottomrule
\end{tabular}
\end{table}

\begin{remark}[All four qualify as Theorems]
All four results are self-contained, complete, and gap-free.  None of the
proofs reduces to a pure definition-unwrap: each requires a structural
observation beyond the definitions.
\begin{itemize}
  \item T35: non-mergeability in $\Fam_{16}$ enforces distinct nodes.
  \item T36: the binary-tree identity $k$ nodes $\Rightarrow$ $\leq k$
        intermediates is a structural tree-counting fact.
  \item T37: distinct transcendental occurrences require distinct nodes
        because no single $\Fam_{16}$ operator can absorb two independent
        transcendentals.
  \item T40: induction with an exact cost-per-step established via T38
        and T40-ADD.
\end{itemize}
\end{remark}

\begin{remark}[T35 dual nature]
The P1 (non-negativity) variant of T35 is admittedly definitional (Cost is
a cardinality, cardinalities are non-negative).  The COST-3 statement
$\Cost(E) \geq |\mathcal{O}(E)|$ is the structurally interesting form and
unambiguously qualifies as a Theorem.  If the document retains P1 as a
separate labelled result, it should be kept as a Proposition or Corollary
of the definition; the T35 label should point to the stronger COST-3
statement.
\end{remark}

\begin{remark}[Alignment with Cost\_Theory\_Complete.tex]
The theorem index in \texttt{Cost\_Theory\_Complete.tex} already classifies
T35, T36, T37 as THEOREM and T40 (there labelled T42) as THEOREM.  This
audit confirms those classifications are correct and provides the detailed
proof-by-proof justification absent in the index.
\end{remark}

% ── Acknowledgments ──────────────────────────────────────────────────────────
\section*{Acknowledgments}

This document is part of the monogate cost theory series (R18, 2026-04-20).
It constitutes the formal reclassification audit for four results previously
labelled Proposition.  All four are reclassified as Theorem.

\bigskip
\noindent\textit{Almaguer, A.R.\ (2026).
Reclassification Audit: T35, T36, T37, T40 (R18).
monogate research. \url{https://monogate.org}}

\end{document}
